\documentclass[10pt]{article}
\usepackage[T1]{fontenc}
\usepackage[utf8]{inputenc}
\usepackage[french]{babel}
\usepackage{amsmath,amssymb,mathrsfs,mathpazo}
\usepackage{geometry}
\geometry{a4paper,margin=2cm}
\usepackage{graphicx,array,tabularx,tabularray,colortbl,multirow,multicol,booktabs,enumitem,pifont,ulem,fancybox,diagbox,cancel}
\usepackage{pstricks,pstricks-add,pst-plot,pst-tree,pst-node,pst-text,pst-func,pst-math,pst-bspline,pst-3dplot,pst-eucl,pst-solides3d}
\let\pscalculate\undefined
\usepackage{pst-bezier}
\usepackage{tikz,pgfplots,tkz-tab}
\usetikzlibrary{arrows,arrows.meta,patterns,calc,positioning,trees,decorations.markings,shapes.geometric,shapes.misc,shapes.arrows,intersections,decorations.pathmorphing,decorations.pathreplacing,backgrounds,shadows,babel}
\pgfplotsset{compat=1.18}
\usepgfplotslibrary{fillbetween}
\usepackage[np]{numprint}
\usepackage{listings,makecell,pas-tableur,twoopt}
\let\barre\undefined
\usepackage[pstricks]{tablvar}
\definecolor{darkgreen}{RGB}{0,100,0}
\providecommand{\red}{\color{red}}
\providecommand{\blue}{\color{blue}}
\providecommand{\green}{\color{green}}
\providecommand{\black}{\color{black}}
\providecommand{\white}{\color{white}}
\providecommand{\gray}{\color{gray}}
\providecommand{\cyan}{\color{cyan}}
\providecommand{\magenta}{\color{magenta}}
\providecommand{\yellow}{\color{yellow}}
\providecommand{\violet}{\color{violet}}
\providecommand{\orange}{\color{orange}}
\providecommand{\euro}{\texteuro}
\providecommand{\EUR}{\texteuro}
\providecommand{\ieme}{\textsuperscript{e}}
\providecommand{\degres}{\ensuremath{^\circ}}
\providecommand{\degree}{\ensuremath{^\circ}}
\providecommand{\vv}[1]{\overrightarrow{#1}}
\providecommand{\Coord}[1]{#1}
\providecommand{\decosix}{}
\providecommand{\starredbullet}{\ensuremath{\bullet}}
\pagestyle{empty}
\setlength{\parindent}{0pt}
\usepackage{adjustbox}\DeclareUnicodeCharacter{2236}{:}\DeclareUnicodeCharacter{2190}{\ensuremath{\leftarrow}}
\providecommand{\R}{\mathbb{R}}
\providecommand{\N}{\mathbb{N}}
\providecommand{\Z}{\mathbb{Z}}
\providecommand{\Q}{\mathbb{Q}}
\providecommand{\D}{\mathbb{D}}
\providecommand{\C}{\mathbb{C}}
\providecommand{\e}{{\mathrm{e}}}
\providecommand{\vect}[1]{\overrightarrow{#1}}
\providecommand{\vectt}[1]{\overrightarrow{\mathrm{#1}}}
\providecommand{\vvt}[1]{\overrightarrow{\mathrm{#1}}}
\providecommand{\np}[2][]{{#2~#1}}
\providecommand{\vv}[1]{{\overrightarrow{#1}}}
\renewcommand{\d}{{\mathrm{d}}}
\providecommand{\degres}{{^\circ}}
\providecommand{\degre}{{^\circ}}
\providecommand{\blue}{\color{blue}}
\providecommand{\red}{\color{red}}
\providecommand{\enskip}{\quad}
\providecommand{\llbracket}{\lbrack\!\lbrack}
\providecommand{\rrbracket}{\rbrack\!\rbrack}
\providecommand{\euro}{{\text{€}}}
\providecommand{\no}{{\text{n°}}}
\providecommand{\interval}[2]{{[#1;#2]}}
\providecommand{\intervFO}[2]{{[#1;#2[}}
\providecommand{\dotfill}{\ldots}

\providecommand{\ds}{\displaystyle}
\providecommand{\textding}[1]{#1}
\providecommand{\Oij}{\(\mathrm{O};\vec{\imath},\vec{\jmath}\)}
\providecommand{\Oijk}{\(\mathrm{O};\vec{\imath},\vec{\jmath},\vec{k}\)}
\providecommand{\Ouv}{\(\mathrm{O};\vec{u},\vec{v}\)}
\providecommand{\intervFF}[2]{\ensuremath{[#1;#2]}}
\renewcommand{\euro}{\ifmmode\text{\texteuro}\else\texteuro\fi}
\DeclareUnicodeCharacter{00B0}{\ifmmode^{\circ}\else\textdegree\fi}
\DeclareUnicodeCharacter{20AC}{\ifmmode\text{\texteuro}\else\texteuro\fi}
% Énoncés : Éducation nationale. Transcription et corrigés : APMEP, Denis Vergès et contributeurs. Découpage et classement : CoopMaths, Sébastien Lozano, Éric Elter et contributeurs. Adaptation de la présentation : Maths Explorer.
% CC BY-SA : https://coopmaths.fr/www/about/
% https://coopmaths.fr/alea/static/bac/2023/tex/bac_2023_09_sujet2_ameriquesud_2_cor.tex
\begin{document}
\noindent\textit{APMEP / CoopMaths. CC BY-SA. Présentation adaptée par Maths Explorer.}\par\medskip
\medskip

Dans un repère orthonormé \Oijk{} on considère les points :
\begin{center} A$(1~;~1~;~-4)$,\quad B$(2~;~-1~;~-3)$,\quad  C$(0~;~-1~;~-1)$\quad  et\quad  $\Omega$(1~;~1~;~2).\end{center}

\smallskip

\begin{enumerate}
\item
Le vecteur $\vectt{AB}$ a pour coordonnées
$\begin{pmatrix}
2-1 \\ -1-1 \\-3-(-4)
\end{pmatrix}
=
\begin{pmatrix}
1 \\ -2 \\ 1
\end{pmatrix}$
 et le vecteur $\vectt{AC}$ a pour coordonnées
$\begin{pmatrix}
0-1 \\ -1-1 \\-1-(-4)
\end{pmatrix}
=
\begin{pmatrix}
-1 \\ -2 \\ 3
\end{pmatrix}$.

Ces deux vecteurs ne sont pas colinéaires, donc les points A, B et C ne sont pas alignés; ils définissent donc un plan.
 
\item 
	\begin{enumerate}
		\item Soit  $\vect{n}$ le vecteur de coordonnées $\begin{pmatrix}1\\1\\1\end{pmatrix}$.
		
\begin{list}{\textbullet}{}
\item $\vect{n}\cdot \vectt{AB}=1\times 1+1\times (-2) +1\times 1=0$ donc $\vect{n}\perp \vectt{AB}$;
\item $\vect{n}\cdot \vectt{AC}=1\times (-1)+1\times (-2) +1\times 3=0$ donc $\vect{n}\perp \vectt{AC}$.
\end{list}		

Le vecteur $\vect{n}$ est donc orthogonal à deux vecteurs directeurs du plan (ABC), donc c'est un vecteur normal au plan (ABC).
		
		\item 
Le plan (ABC) a donc une équation cartésienne de la forme $x+y+z+d=0$. On détermine $d$ en utilisant le fait que le point A appartient au plan (ABC).

$\text{A}\in (\text{ABC})
\iff
x_{\text A} + y_{\text A} + z_{\text A} +d=0
\iff
1+1-4+d=0
\iff d=2$

Le plan (ABC) a donc pour équation cartésienne $x+y+z+2=0$.
	
	\end{enumerate}
	
\item 
	\begin{enumerate}
		\item 
$x_{\Omega}	 + y_{\Omega} + z_{\Omega} +2 = 1+1-2+2=2\neq 0$ donc $\Omega \notin (\text{ABC})$.
		
		\item Soit H le projeté orthogonal du point $\Omega$ sur le plan (ABC).
		
Donc $\vect{\Omega H}$ est un vecteur orthogonal au plan (ABC), donc il est colinéaire au vecteur $\vect n$, et il a donc des coordonnées de la forme
$\begin{pmatrix}
k \\k \\k
\end{pmatrix}$, où $k\in\R$.

Les coordonnées de $\vect{\Omega H}$ sont
$\begin{pmatrix}
x_{\text H} - x_{\Omega} \\ y_{\text H} - y_{\Omega} \\ z_{\text H} - z_{\Omega}
\end{pmatrix}
=
\begin{pmatrix}
x_{\text H} - 1 \\ y_{\text H} -1 \\ z_{\text H} -2
\end{pmatrix}$.

On déduit donc que
$\left \{
\begin{array}{l !{=} l}
x_{\text H} -1& k\\
y_{\text H} -1& k\\
z_{\text H} -2& k 
\end{array}
\right .$, c'est-à-dire
$\left \{
\begin{array}{l !{=} r}
x_{\text H} & 1+k\\
y_{\text H} & 1+k\\
z_{\text H} & 2+k 
\end{array}
\right .$

Or le point H appartient au plan (ABC) donc:
$x_{\text H} + y_{\text H} + z_{\text H} +2=0
\iff
(1+k)+(1+k)+(2+k)+2=0
\iff
6+3k=0
\iff
k=-2$.

Le point H a donc pour coordonnées
$\begin{pmatrix}
1-2 \\ 1-2 \\ 2-2
\end{pmatrix}$
soit
$\begin{pmatrix}
-1 \\ -1 \\ 0
\end{pmatrix}$.
	\end{enumerate}
			\end{enumerate}

On admet que $\Omega \text{H} = 2\sqrt 3$.
On définit la sphère $S$ de centre $\Omega$ et de rayon $2\sqrt 3$ comme l'ensemble de tous les points M de l'espace tels que $\Omega\text{M} = 2\sqrt 3$.

	\begin{enumerate}[resume]

\item 
Le point H est le projeté orthogonal de $\Omega$ sur le plan (ABC), donc $\Omega$H est la plus courte distance entre $\Omega$ et tout point du plan; autrement dit, tout point N de (ABC) distinct de H est tel que $\text{$\Omega$N} > \text{$\Omega$H}$, donc $\text{$\Omega$N} > 2\sqrt{3}$.

Ce qui prouve que le point N n'appartient pas à la sphère $S$.
\end{enumerate}

On dit qu'un plan $\mathcal{P}$ est tangent à la sphère $S$ en un point K lorsque les deux conditions suivantes sont vérifiées :

\hfill$\begin{array}{l}
\bullet~~ \text{K} \in \mathcal{P} \cap S\\
\bullet~~ (\Omega\text{K}) \perp \mathcal{P}
\end{array}$\hfill\null

	\begin{enumerate}[resume]

\item Soit le plan $\mathcal{P}$ d'équation  $x +y - z - 6 = 0$ et le point K de coordonnées K(3~;~3~;~0).



\begin{list}{\textbullet}{}
\item $x_{\text K} + y_{\text K} - z_{\text K} -6 = 3 + 3 - 0 - 6 = 0$ donc $\text K \in \mathcal{P}$

$\text{$\Omega$K}^2 = \left (3-1\right )^2 + \left ( 3-1\right )^2 + \left (0-2\right )^2 = 4+4+4=12$ donc $\text{$\Omega$K} =2\sqrt{3}$; et donc $\text K \in S$.

Donc $\text{K} \in \mathcal{P} \cap S$.
\item Le vecteur $\vect{\Omega K}$ a pour coordonnées 
$\begin{pmatrix}
3-1 \\ 3-1\\ 0-2
\end{pmatrix}
=
\begin{pmatrix}
2 \\ 2\\ -2
\end{pmatrix}$.

Le plan $\mathcal{P}$ a pour vecteur normal 
$\vect{n'} \begin{pmatrix}
1 \\ 1\\ -1
\end{pmatrix}$

$\vect{\Omega K} = -2 \vect{n'}$; or $\vect{n'}$ est un vecteur normal au plan $\mathcal{P}$ donc $\vect{\Omega K}$ est orthogonal au plan $\mathcal{P}$, et donc $(\text{$\Omega$K}) \perp \mathcal{P}$.
\end{list}

On en conclut que  le plan $\mathcal{P}$ est tangent à la sphère $S$ au point K.

\item On admet que les plans (ABC) et $\mathcal{P}$ sont sécants selon une droite $(\Delta)$.

\emph{Rem. } -- Le plan (ABC) a pour vecteur normal $\vect{n}\begin{pmatrix}1\\1\\1\end{pmatrix}$  et $\mathcal{P}$ a pour vecteur normal $\vect{n'}\begin{pmatrix}1\\1\\-1\end{pmatrix}$.

Ces vecteurs ne sont manifestement pas colinéaires donc (ABC) et $\mathcal{P}$ sont sécants selon une droite $(\Delta)$.

\smallskip

On résout le système :

$M(x~;~y~;~z) \in (\text{ABC}) \:\cap \:\mathcal{P} \iff \left \{
\begin{array}{r}
x+y+z+2=0\\
x+y-z-6=0
\end{array}
\right.
\iff $

$\left\{
\begin{array}{r}
x+y+z+2=0\\
(x+y+z+2)-(x+y-z-6)=0
\end{array}
\right.\iff
\left\{
\begin{array}{r}
x+y+z+2=0\\
z+2+z+6=0
\end{array}
\right .
\iff
\left \{
\begin{array}{r !{=} r}
x+y+z+2 & 0\\
z&-4
\end{array}
\right .$

$\iff
\left \{
\begin{array}{r !{=} r}
x + y -2& 0\\
z&-4
\end{array}
\right .$

On peut l'écrire
$\left \{
\begin{array}{r !{=} r}
x& 2-y\\
y&y\\
z&-4\phantom{-y}
\end{array}
\right .$
ou encore
$\left \{
\begin{array}{r !{=} r}
x& 2-t\\
y&t\\
z&-4\phantom{-t}
\end{array}
\right .$ où $t\in\R$.

La droite $(\Delta)$ a donc pour représentation paramétrique:
$\left \{
\begin{array}{r !{=} r}
x& 2-t\\
y&t\\
z&-4\phantom{-t}
\end{array}
\right .$ où $t\in\R$.
\end{enumerate}

\end{document}
